\section{Fractional spectra and the limits of singularity degree}
\label{sec:fractional-sdp}

Semidefinite sublevels need not have either of the polyhedral diameter
rates. We analyze a trace-normalized version of the $3\times3$ instance
in \citet[Example 2]{sturm2000}. The associated nested constraints and
quarter-power error geometry are classical; see also
\citet[Theorem 4.5.1, Example 4.5.2, and Section 4.7]{drusvyatskiy2017}.
Our purpose is to compute its entire reduced log-barrier spectrum,
transfer the resulting scales to every fixed barrier, and compare it
with a section having the same singularity degree but a different
conditioning exponent.

\subsection{Two spectrahedra and their sharp diameters}

Consider
\begin{equation}
 \min X_{22}\quad\text{subject to}\quad
 X\in\mathbb S^3_+,\quad \tr X=1,\quad X_{33}=X_{12},
 \label{eq:fractional-sdp}
\end{equation}
and the restricted problem obtained by adding
\begin{equation}
 X_{13}=X_{23}=0.
 \label{eq:restricted-sdp}
\end{equation}
Here $\mathbb S^3_+$ is the cone of real symmetric positive
semidefinite matrices. Write their feasible sets as $P$ and $P_0$,
and their sublevel diameters as $D$ and $D_0$, respectively, each
in its inherited Frobenius metric.

\begin{proposition}[Equal singularity degree, different conditioning]
\label{prop:paired-sdp}
Both problems have compact feasible sets with positive-definite
feasible points and unique optimum $X^*=e_1e_1^\top$. The optimality
systems of both problems have singularity degree two. Nevertheless,
\begin{equation}
 D(g)=\Theta(g^{1/4}),\qquad D_0(g)=\Theta(g^{1/2}).
 \label{eq:paired-diameters}
\end{equation}
Thus every fixed self-concordant barrier on the respective set satisfies
\begin{equation}
 \widehat\kappa_F(g)=\Theta(g^{-3/2})\quad\text{on }P,
 \qquad
 \widehat\kappa_{F_0}(g)=\Theta(g^{-1})\quad\text{on }P_0.
 \label{eq:paired-conditioning}
\end{equation}
For the displayed conic formulations, neither problem has a strictly
complementary optimal primal--dual pair.
\end{proposition}

We use the standard definition of singularity degree for a feasible
conic \emph{system}: it is the minimum number of facial-reduction
steps required to reach a face in whose relative interior the affine
system has a feasible point \citep{drusvyatskiy2017,wakimuramatsu2013}.
At each step an equality combination nonnegative on the current cone
and zero on the affine system exposes a smaller face. In the proposition
the system includes the optimality equality $X_{22}=0$.
The original feasible systems have singularity degree zero, because
they contain positive-definite matrices. Keeping these systems distinct
is essential when applying error bounds to objective accuracy.

\begin{proof}
Trace normalization gives compactness. The matrix
\[
 \begin{pmatrix}2/5&1/5&0\\1/5&2/5&0\\0&0&1/5\end{pmatrix}
\]
is positive definite and feasible for both problems. Write a point of
$P$ as
\begin{equation}
 X=\begin{pmatrix}a&b&t\\b&u&e\\t&e&b\end{pmatrix},
 \qquad a+u+b=1.
 \label{eq:fractional-coordinates}
\end{equation}
If $u=0$, the second row and column of a PSD matrix vanish, so
$b=e=0$. The third diagonal is then zero, so $t=0$. Thus the
unique optimum in both sets is $X^*$, with value zero.

For $X\in L(g)$, principal minors give
\begin{equation}
 0\leq b\leq\sqrt g,\qquad
 t^2\leq ab\leq\sqrt g,\qquad
 e^2\leq ub\leq g^{3/2},\qquad
 |a-1|=u+b\leq g+\sqrt g.
 \label{eq:fractional-minors}
\end{equation}
Together with $0\leq u\leq g$, these inequalities prove
$\norm{X-X^*}_F=O(g^{1/4})$. If $t=e=0$, they instead give
$\norm{X-X^*}_F=O(g^{1/2})$. These prove both diameter upper
bounds.

For $0<g\leq1/16$, set
\begin{equation}
 u=g,\qquad b=\frac{\sqrt g}{2},\qquad
 a=1-g-\frac{\sqrt g}{2},\qquad e=0,\qquad
 t=\pm\frac{g^{1/4}}{2\sqrt2}.
 \label{eq:fractional-chords}
\end{equation}
Here $a\geq13/16$ and $ag-b^2=g(a-1/4)>0$, while
\[
 \det X=agb-b^3-gt^2
         =g^{3/2}(a/2-1/4)>0.
\]
The leading principal minors are positive, so both matrices are
positive definite. Their Frobenius distance is $g^{1/4}$, proving
the first lower bound. For $P_0$, take the same $a,b,u$ and
$t=e=0$. The resulting matrix is positive definite and its distance
from $X^*$ is at least $\sqrt2b=\sqrt{g/2}$, proving the second.

For the facial-reduction claim let
$E_{ii}=e_ie_i^\top$ and
$J_{ij}=(e_ie_j^\top+e_je_i^\top)/2$, so that
$\inner{J_{ij}}{X}=X_{ij}$. The first equality $X_{22}=0$
exposes the face supported on $\operatorname{span}\{e_1,e_3\}$.
On this face $X_{12}=0$, so $X_{33}-X_{12}=0$ exposes the ray
generated by $E_{11}$. Its relative interior contains $X^*$, proving
that two steps suffice in both systems.

To exclude one step, an exposing combination for either optimality
system has the form
\begin{equation}
 Z=\alpha I+\beta(E_{33}-J_{12})+\gamma E_{22}
       +\delta J_{13}+\eta J_{23}\succeq0,
 \label{eq:exposing-combination}
\end{equation}
where $\delta=\eta=0$ for $P$. Its value on the affine system
must be zero; the trace equation is the only equation with nonzero
right side, so $\alpha=0$. Equivalently $Z_{11}=0$.
A PSD matrix with zero diagonal entry has a zero corresponding row.
The first row therefore forces $\beta=\delta=0$; then
$Z_{33}=0$ forces $\eta=0$. Consequently every nonzero first-step
certificate is a positive multiple of $E_{22}$. It cannot expose
the final ray. This proves minimality of the two-step chain.

The same calculation establishes the complementarity assertion.
A dual slack has the form \eqref{eq:exposing-combination} with
$\gamma=1$, since the objective is $X_{22}$; optimality makes
$\alpha=0$. The only optimal slack is therefore $E_{22}$.
It has rank one, as does $X^*$, so their ranks sum to two rather
than three. Dual strict feasibility nonetheless holds: choosing a
negative multiplier for the trace equation and all other multipliers
zero makes the slack $E_{22}+\tau I\succ0$ for $\tau>0$.
Finally, \cref{thm:diameter-law} proves
\eqref{eq:paired-conditioning} for every fixed barrier.
\end{proof}

For a singleton SDP optimum, the classical singularity-degree error
bound yields $D(g)=O(g^{2^{-s}})$ when the optimality system has
degree $s$, on a fixed compact set \citep{sturm2000,drusvyatskiy2017}.
Indeed, a feasible point has zero cone residual and violates only the
added optimality equation by $g$; its distance to the optimal affine
space is at most a fixed multiple of that residual, by a right inverse
on the range of the equality map. The error bound then controls its
distance to the optimum. If this exponent is attained, the diameter
law gives
\begin{equation}
 \widehat\kappa_F(g)=\Theta(g^{\,2^{1-s}-2}).
 \label{eq:attained-singularity-rate}
\end{equation}
The two degree-two systems above show why attainment must be verified:
the first has the quarter-power diameter, whereas the second has a
strictly faster square-root decay. Thus singularity degree alone does
not determine the conditioning exponent, even for these compact
$3\times3$ formulations. The diameter records sublevel shrinkage
that this facial-reduction count does not determine.

\subsection{The exact log-barrier center and four spectral scales}

The general diameter theorem gives the condition-number exponent on
$P$. In this example one can also compute every eigenvalue scale.
This supplements the geometric argument with an exact spectral
description, including its leading constants.

\begin{theorem}[Four reduced-Hessian scales]
\label{thm:fractional-spectrum}
For $F(X)=-\log\det X$ on \eqref{eq:fractional-sdp}, the central
point at gap $g$ has the form
\begin{equation}
 \widehat X_{\log}(g)=
 \begin{pmatrix}a&b&0\\b&g&0\\0&0&b\end{pmatrix},\qquad
 b=\frac{-g+\sqrt{3g-2g^2}}{3},\quad a=1-g-b,
 \label{eq:fractional-exact-center}
\end{equation}
and the corresponding central-path parameter satisfies
\begin{equation}
 q=ag-b^2,\qquad
 \mu=\frac{q}{1-b-2g}.
 \label{eq:fractional-gap-mu}
\end{equation}
These formulas apply for all sufficiently small positive $g$.
For the four increasingly ordered eigenvalues of the reduced Hessian
in the Frobenius metric, as $\mu\downarrow0$,
\begin{align}
 \lambda_1&\sim\sqrt2\,\mu^{-1/2},&
 \lambda_2&\sim\mu^{-1},&
 \lambda_3&\sim\sqrt2\,\mu^{-3/2},&
 \lambda_4&\sim\frac47\mu^{-2}.
 \label{eq:four-eigenvalues}
\end{align}
In particular,
\begin{equation}
 \kappa_{\log}(X_{\log}(\mu))
 \sim\frac{2\sqrt2}{7}\mu^{-3/2}
 \sim\frac{3\sqrt3}{7}g^{-3/2}.
 \label{eq:fractional-condition-constant}
\end{equation}
On $P_0$, the log-barrier center is the same matrix. Its two
reduced-Hessian eigenvalues are asymptotic to
$\mu^{-1}$ and $(4/7)\mu^{-2}$, respectively, so its condition
number is asymptotic to $(4/7)\mu^{-1}$, or $(6/7)g^{-1}$.
\end{theorem}

\begin{proof}
Conjugation by $\Diag(1,1,-1)$ preserves the feasible set,
the objective, and $F$. Uniqueness of the central point therefore
forces $t=e=0$ in \eqref{eq:fractional-coordinates}. At a fixed
gap $u=g$, the remaining free coordinate is $b$, and
\[
 F=-\log b-\log q,\qquad
 q=g(1-g-b)-b^2.
\]
Stationarity in $b$ gives $q+bq_b=0$, or
\[
 3b^2+2gb+g^2-g=0.
\]
The positive root is the value in
\eqref{eq:fractional-exact-center}. Stationarity of $g/\mu+F$ in $g$ gives
$1/\mu=q_g/q$, which is \eqref{eq:fractional-gap-mu}.
For small $g>0$, $b>0$, $a>0$, $q>0$, and $q_g>0$.
Thus these matrices are interior stationary points for a positive
$\mu$, hence are the unique central points by strict convexity.
Conversely uniqueness and the fixed-gap stationarity determine every
small-gap center. The same argument holds on $P_0$, where $t=e=0$
are imposed as constraints.

For an explicit parametrization of the entire small-parameter tail, put
$\tau=g/b$ and $d=\tau^2+2\tau+3$. The stationarity equations become
\begin{equation}
 a=\frac{\tau+3}{d},\qquad b=\frac\tau d,\qquad
 g=\frac{\tau^2}{d},\qquad q=\frac{\tau^2(\tau+2)}{d^2},\qquad
 \mu=\frac{\tau^2(\tau+2)}{d(3+\tau-\tau^2)}.
 \label{eq:fractional-rational-parameter}
\end{equation}
For $0<\tau<1$, all required positivity conditions hold. Both $g$ and
$\mu$ extend continuously to $[0,1]$, taking values $0$ and $1/6$ at
the endpoints. Thus every gap and every central parameter in $(0,1/6)$
is attained. Moreover $\tau^2\leq6g$ and $\tau^2\leq12\mu$ on this
interval, so any such parameter choices tend to zero as $g$ or $\mu$
tends to zero. Consequently the limits obtained from
\eqref{eq:fractional-rational-parameter} describe the entire central
path tail, rather than a selected sequence of central points.

As $g\downarrow0$, the formulas give
\begin{equation}
 b\sim\sqrt{g/3},\quad q\sim2g/3,\quad
 \mu\sim2g/3;\qquad
 a\to1,\quad b\sim\sqrt{\mu/2},\quad q\sim\mu.
 \label{eq:fractional-basic-asymptotics}
\end{equation}
We now compute the Hessian in the four tangent coordinates
$(b,g,t,e)$. The $(b,g)$ and $(t,e)$ subspaces are orthogonal
in the Frobenius metric, and the Hessian has no mixed terms between
them at $t=e=0$, either by the same sign symmetry or by direct
differentiation.

Let $A_2=\begin{pmatrix}a&b\\b&g\end{pmatrix}$, with determinant
$q$. For an off-block tangent perturbation
$U=\begin{pmatrix}0_{2\times2}&v\\v^\top&0\end{pmatrix}$,
where $v=(t,e)^\top$, one has
\[
 D^2F(X)[U,U]=\tr(X^{-1}UX^{-1}U)
             =\frac2b v^\top A_2^{-1}v,
 \qquad \norm{U}_F^2=2v^\top v.
\]
Thus this block's eigenvalues are those of $A_2^{-1}/b$.
The larger eigenvalue of $A_2$ tends to one, and the smaller is
asymptotic to $q\sim\mu$, by their product. This gives precisely
$\sqrt2\mu^{-1/2}$ and $\sqrt2\mu^{-3/2}$.

On the $(b,g)$ subspace, the tangent Gram matrix and the coordinate
Hessian $K$ are
\begin{equation}
 G=\begin{pmatrix}4&1\\1&2\end{pmatrix},\qquad
 K=\begin{pmatrix}
 b^{-2}+q_b^2/q^2+2/q&q_bq_g/q^2+1/q\\
 q_bq_g/q^2+1/q&q_g^2/q^2+2/q
 \end{pmatrix},
 \label{eq:fractional-generalized-block}
\end{equation}
where $q_b=-g-2b$ and $q_g=1-b-2g$. In particular the
coordinate Hessian $K$ alone does not give the Euclidean operator's
eigenvalues; they solve $Kv=\lambda Gv$. From
\eqref{eq:fractional-basic-asymptotics},
\begin{equation}
 \mu K_{bb}\to6,\qquad
 \mu^{3/2}K_{bg}\to-\sqrt2,\qquad
 \mu^2K_{gg}\to1.
 \label{eq:fractional-block-limits}
\end{equation}
It follows that
$\mu^2G^{-1/2}KG^{-1/2}$ converges to the rank-one matrix
$G^{-1/2}e_2e_2^\top G^{-1/2}$, whose nonzero eigenvalue is
$(G^{-1})_{22}=4/7$. Hence the larger generalized eigenvalue
is asymptotic to $(4/7)\mu^{-2}$. Also
\[
 \mu^3\det K\longrightarrow6-2=4,\qquad \det G=7.
\]
The product of the generalized eigenvalues is $\det K/\det G$,
so the smaller is asymptotic to $\mu^{-1}$. The four distinct
powers determine their eventual increasing order, proving
\eqref{eq:four-eigenvalues}. Their extreme ratio proves
\eqref{eq:fractional-condition-constant}, using
$g\sim3\mu/2$. Restricting to $P_0$ retains precisely the
$(b,g)$ block and gives its two stated eigenvalues and ratio.
\end{proof}

The log-barrier calculation in \cref{thm:fractional-spectrum} has a
Lean formalization in the accompanying repository. It proves that the
matrix is the unique minimizer of the actual barrier objective over all
positive-definite feasible matrices, computes the second Fr\'echet
derivative of $-\log\det$, identifies its operator in the Frobenius metric,
and factors its characteristic polynomial. The formalization also proves
the ordered asymptotics and both condition-number constants, including the
restricted problem and coverage of the full small-parameter tail. The
verification report is \texttt{formal/FRACTIONAL\_SDP.md}. Its scope is the
log-barrier spectrum; the singularity-degree and diameter assertions in
\cref{prop:paired-sdp}, and the following transfer to other barriers, are
separate mathematical results outside this formalization.

At equal gap the Loewner comparison in
\cref{thm:difference-body} transfers more than the extreme ratio.
Every fixed barrier on $P$ has ordered eigenvalues of orders
$g^{-1/2}$, $g^{-1}$, $g^{-3/2}$, and $g^{-2}$ in the same
Frobenius metric. Every fixed barrier on $P_0$ has orders
$g^{-1}$ and $g^{-2}$. The leading constants in
\eqref{eq:four-eigenvalues} concern the canonical log barrier;
they are not asserted to be invariant under barrier changes.

SDP spectral clustering has substantial precedent. In particular,
\citet[Section 4]{alizadeh1998} analyze eigenvalue behavior of a
primal--dual Schur matrix, and \citet[Section 3.3]{pena2002}
compares central-path system conditioning. Those results concern a
different matrix and, in the cited SDP clustering setting, strict
complementarity. \citet{wei2010} construct instances with prescribed
failure of strict complementarity and investigate their numerical
difficulty.

\citet[Section 4]{sremac2021} relate singularity degree to spectral rates
along an external regularization path: $X(\alpha)$ maximizes
$\log\det X$ subject to $A(X)=b+\alpha A(B)$, with fixed $B\succ0$,
and $Z(\alpha)=\alpha X(\alpha)^{-1}$. Under their assumptions,
Corollary 4.3 separates eigenvalue rates of $X(\alpha)$ and
$Z(\alpha)$, Theorem 4.4 characterizes degree one through linear
decay of the vanishing primal eigenvalues, and Theorem 4.7 gives
separated rates for dual diagonal blocks. These are relevant
antecedents for the degree--spectrum connection. Here the affine
constraints remain fixed, the parameter is the primal objective gap,
and the operator is the equality-reduced Hessian. Primal-iterate
eigenvalues alone do not determine its compression to the equality
tangent space.

The calculation above concerns the equality-reduced
primal log-det Hessian at a non-strictly-complementary optimum, with
its external metric specified. To the best of our knowledge, the four
leading constants in \eqref{eq:four-eigenvalues} have not previously
been stated for this trace-normalized classical construction.
The novelty claimed here is this explicit spectral calculation and
its all-barrier transfer, not the construction of a fractional SDP
error bound or the general study of SDP spectral clusters.
